% GENERATED FILE. Edit canonical lessons, then run npm run generate:books.
% canonical-source-hash: fnv1a64:451df18e8633e7a7
% canonical-lessons: PT-C326-lenda, PT-C326-lendario, PT-C326-mito, PT-C326-fabula, PT-C326-fada

\chapter{Lendas e mitos}
\label{ch:pt-lendas-e-mitos}
\hlchaptermodality{326}

\begin{chapteropening}
\textbf{By the end of this chapter:} \emph{I can talk about legends, myths, fables and fairies.}

Complete the last lesson of chapter 326: I can talk about legends, myths, fables and fairies.
\end{chapteropening}

\section[a lenda]{a lenda — the legend}

\label{lesson:PT-C326-lenda}



\begin{quote}
Before the new one: say the Portuguese for the villain, then the Portuguese for the ending.
\end{quote}

\subsection*{You'll want to know: a lenda}
\textbf{a lenda} — \textquotedblleft{}the legend\textquotedblright{}.

\begin{quote}
Segundo a lenda, ninguém pode entrar na gruta à noite.
\end{quote}

\begin{quote}
\emph{According to the legend, nobody may go into the cave at night.}
\end{quote}

\emph{Conta a lenda que…} — legend has it that…

\begin{cousinweb}
\emph{lenda} comes down from Old Portuguese \emph{leenda} (EtymDB), from Medieval Latin \emph{legenda} \textquotedblleft{}things to be read\textquotedblright{}, from \emph{legere} \textquotedblleft{}to read\textquotedblright{}.
\end{cousinweb}

\begin{grammarlens}[title={What you've built}]
One more word for telling, explaining and describing at B1.
\end{grammarlens}

\subsection*{Your turn}
\begin{itemize}
\raggedright
  \item \emph{Say it:} \emph{a lenda}
  \item \emph{Say it:} \emph{a lenda}, once more
  \item \emph{Say it:} a sentence of your own with \emph{a lenda}
  \item \emph{Recall:} say the Portuguese for the villain, then the Portuguese for the ending, then say \emph{a lenda} again
\end{itemize}

\begin{tcolorbox}[breakable,colback=teal!4,colframe=teal!35!black,title={Before you move on}]
What does a lenda mean? (\textquotedblleft{}The legend\textquotedblright{}.) Say it once more.
\end{tcolorbox}
\section[lendário]{lendário — legendary}

\label{lesson:PT-C326-lendario}



\begin{quote}
Before the new one: say the Portuguese for the ending, then the Portuguese for the legend.
\end{quote}

\subsection*{You'll want to know: lendário}
\textbf{lendário} — \textquotedblleft{}legendary\textquotedblright{}.

\begin{quote}
Aquele cantor é lendário na região.
\end{quote}

\begin{quote}
\emph{That singer is legendary in the region.}
\end{quote}

Built on \emph{lenda}.

\begin{cousinweb}
\emph{lendário} is built on \emph{lenda} \textquotedblleft{}a legend\textquotedblright{}, from Medieval Latin \emph{legenda} (EtymDB), with the ending \emph{-ário}.
\end{cousinweb}

\begin{grammarlens}[title={What you've built}]
One more word for telling, explaining and describing at B1.
\end{grammarlens}

\subsection*{Your turn}
\begin{itemize}
\raggedright
  \item \emph{Say it:} \emph{lendário}
  \item \emph{Say it:} \emph{lendário}, once more
  \item \emph{Say it:} a sentence of your own with \emph{lendário}
  \item \emph{Recall:} say the Portuguese for the ending, then the Portuguese for the legend, then say \emph{lendário} again
\end{itemize}

\begin{tcolorbox}[breakable,colback=teal!4,colframe=teal!35!black,title={Before you move on}]
What does lendário mean? (\textquotedblleft{}Legendary\textquotedblright{}.) Say it once more.
\end{tcolorbox}
\section[o mito]{o mito — the myth}

\label{lesson:PT-C326-mito}



\begin{quote}
Before the new one: say the Portuguese for the legend, then the Portuguese for legendary.
\end{quote}

\subsection*{You'll want to know: o mito}
\textbf{o mito} — \textquotedblleft{}the myth\textquotedblright{}.

\begin{quote}
É um mito: a água fria não faz mal.
\end{quote}

\begin{quote}
\emph{It is a myth: cold water does no harm.}
\end{quote}

Also a false belief, as here.

\begin{cousinweb}
\emph{mito} comes from Greek \emph{mŷthos} \textquotedblleft{}a word, a speech, a story\textquotedblright{} (EtymDB).
\end{cousinweb}

\begin{grammarlens}[title={What you've built}]
One more word for telling, explaining and describing at B1.
\end{grammarlens}

\subsection*{Your turn}
\begin{itemize}
\raggedright
  \item \emph{Say it:} \emph{o mito}
  \item \emph{Say it:} \emph{o mito}, once more
  \item \emph{Say it:} a sentence of your own with \emph{o mito}
  \item \emph{Recall:} say the Portuguese for the legend, then the Portuguese for legendary, then say \emph{o mito} again
\end{itemize}

\begin{tcolorbox}[breakable,colback=teal!4,colframe=teal!35!black,title={Before you move on}]
What does o mito mean? (\textquotedblleft{}The myth\textquotedblright{}.) Say it once more.
\end{tcolorbox}
\section[a fábula]{a fábula — the fable}

\label{lesson:PT-C326-fabula}



\begin{quote}
Before the new one: say the Portuguese for legendary, then the Portuguese for the myth.
\end{quote}

\subsection*{You'll want to know: a fábula}
\textbf{a fábula} — \textquotedblleft{}the fable\textquotedblright{}.

\begin{quote}
Na fábula, o leão e o rato ficam amigos.
\end{quote}

\begin{quote}
\emph{In the fable, the lion and the mouse become friends.}
\end{quote}

A short story with animals and a lesson at the end.

\begin{cousinweb}
\emph{fábula} was borrowed from Latin \emph{fābula} \textquotedblleft{}a story, a tale\textquotedblright{} (EtymDB), from \emph{fārī} \textquotedblleft{}to speak\textquotedblright{}.
\end{cousinweb}

\begin{grammarlens}[title={What you've built}]
One more word for telling, explaining and describing at B1.
\end{grammarlens}

\subsection*{Your turn}
\begin{itemize}
\raggedright
  \item \emph{Say it:} \emph{a fábula}
  \item \emph{Say it:} \emph{a fábula}, once more
  \item \emph{Say it:} a sentence of your own with \emph{a fábula}
  \item \emph{Recall:} say the Portuguese for legendary, then the Portuguese for the myth, then say \emph{a fábula} again
\end{itemize}

\begin{tcolorbox}[breakable,colback=teal!4,colframe=teal!35!black,title={Before you move on}]
What does a fábula mean? (\textquotedblleft{}The fable\textquotedblright{}.) Say it once more.
\end{tcolorbox}
\section[a fada]{a fada — the fairy}

\label{lesson:PT-C326-fada}



\begin{quote}
Before the new one: say the Portuguese for the myth, then the Portuguese for the fable.
\end{quote}

\subsection*{You'll want to know: a fada}
\textbf{a fada} — \textquotedblleft{}the fairy\textquotedblright{}.

\begin{quote}
A fada transformou a abóbora num carro.
\end{quote}

\begin{quote}
\emph{The fairy turned the pumpkin into a carriage.}
\end{quote}

\emph{Conto de fadas} — fairy tale.

\begin{cousinweb}
\emph{fada} goes back to Vulgar Latin \emph{Fāta} \textquotedblleft{}a goddess of fate\textquotedblright{} (EtymDB), from Latin \emph{fātum} \textquotedblleft{}fate\textquotedblright{}, from \emph{fārī} \textquotedblleft{}to speak\textquotedblright{}.
\end{cousinweb}

\begin{grammarlens}[title={What you've built}]
One more word for telling, explaining and describing at B1.
\end{grammarlens}

\subsection*{Your turn}
\begin{itemize}
\raggedright
  \item \emph{Say it:} \emph{a fada}
  \item \emph{Say it:} \emph{a fada}, once more
  \item \emph{Say it:} a sentence of your own with \emph{a fada}
  \item \emph{Recall:} say the Portuguese for the myth, then the Portuguese for the fable, then say \emph{a fada} again
\end{itemize}

\begin{tcolorbox}[breakable,colback=teal!4,colframe=teal!35!black,title={Before you move on}]
What does a fada mean? (\textquotedblleft{}The fairy\textquotedblright{}.) Say it once more.
\end{tcolorbox}
